\chapter{结论与展望}{Conclusion and prospect}

\gdutbacksection{研究结论}
针对立体仓库堆叠货物盘点任务中存在的目标遮挡严重、结构关系复杂、数量推断困难以及模型部署资源受限等问题，本文围绕复杂场景识别、数量推理与部署适配开展研究。以实例分割为基础，结合后处理数量计算方法与模型压缩策略，形成了面向立体仓库堆叠货物盘点任务的技术路线。全文的主要研究结论如下。

\par\noindent
（1）针对复杂堆叠场景下现有实例分割模型易出现漏检、误检以及结构建模能力不足的问题，本文以 YOLO11s-Seg 为基线，提出改进模型 YOLO11-StackLite。该模型通过引入通道拆分注意力模块 C2SASA、重排分组卷积模块 SSConv 以及具有四向卷积结构的 C3k2\_QDConv 模块，从注意力建模、下采样效率和方向性空间表达三个方面对网络结构进行了优化。实验结果显示，改进后的模型平均 box mAP 和平均 mask mAP 分别达到 85.7\% 和 84.7\%，平均 DA 和平均 CA 分别达到 95.88\% 和 94.87\%，表明该结构改进在当前数据集与实验设置下能够提升模型对复杂堆叠场景的适应能力。

\par\noindent
（2）针对仅依赖检测结果难以准确完成货物总数量推断的问题，本文承接前期研究中的面积比后处理计数流程，并完成其与改进模型输出的适配与集成。该流程通过顶层满载检测和基于掩膜关键点的层间面积比估计，对堆叠货物的层次结构进行判别与数量推理；当仅检测到下层目标、视觉证据不足以支撑可靠推理时，则输出 \texttt{num=-1} 并交由系统复核。实验结果显示，在全量验证集上，YOLO11-StackLite 的平均 CA 达到 94.87\%，在顶层未满载有效计数图像这一更困难的子集上，平均 CA 仍达到 80.90\%。该部分工作的贡献边界主要在于既有计数流程的任务适配与实验验证，而不在于提出新的计数方法。

\par\noindent
（3）针对高精度实例分割模型在边缘部署场景中存在参数量较大、计算复杂度较高的问题，本文进一步研究了面向边缘部署需求的模型压缩方法。以 YOLO11-StackLite 为基础，构建了由稀疏约束训练、结构化物理剪枝、保留通道权重迁移、剪枝后微调恢复以及知识蒸馏增强组成的压缩路线，并完成了相应压缩实验验证。实验结果显示，在 30\% 代表性压缩设置下，模型参数量由 9.4 M 降至 7.8 M，GFLOPs 由 34.9 降至 29.1；物理剪枝模型经微调恢复后，平均 DA 和平均 CA 分别达到 94.82\% 和 93.88\%，在进一步蒸馏增强后分别提升至 95.10\% 和 94.02\%。上述结果表明，该压缩路线在当前实验设置下能够在降低模型复杂度的同时较好保持主要任务性能。

综上，本文围绕立体仓库堆叠货物盘点任务，构建了“改进实例分割模型 - 后处理数量估计流程集成 - 模型压缩与部署可行性分析”的研究框架。就论文贡献边界而言，本文的核心在于复杂堆叠场景下实例分割模型的针对性改进、前期计数流程与改进模型输出的有效适配，以及面向边缘部署需求的真实物理剪枝压缩方案验证。其中，数量推理部分的主要价值在于流程衔接与系统验证，而本文的核心原创性仍主要集中在实例分割模型改进和压缩路线设计上。总体上，本文实现的是一种“可判定场景自动输出、不可判定场景保守回退”的智能盘点流程，而非在所有视觉条件下都强制给出闭环自动计数结果。由于实验数据主要来自单一烟草仓库中的规则堆叠场景，且尚未完成实机闭环部署测试，相关结论仍需通过更多跨场景数据和具体设备验证进一步完善。

\gdutbacksection{未来研究展望}
尽管本文围绕立体仓库堆叠货物盘点任务开展了较为系统的研究，并取得了一定成果，但仍存在进一步完善和拓展的空间，主要体现在以下几个方面。

\par\noindent
（1）在数据集方面，本文采用的 TCSD 数据来源于真实仓储场景，能够较好支撑当前研究，但在样本规模、场景多样性以及极端遮挡情形方面仍有进一步扩展空间。后续可进一步采集更多不同仓储条件、不同光照环境、不同货物类型以及不同堆叠方式下的数据，并补充跨仓库、跨场景验证，以增强模型的泛化能力和鲁棒性。

\par\noindent
（2）在数量计算方法方面，本文所集成的后处理算法主要基于实例分割掩膜与几何先验进行结构推理。对于更复杂的不规则堆叠、严重遮挡或局部缺失场景，现有方法仍可能受到一定影响。后续可考虑引入更强的结构约束、三维几何信息或多视角信息，以进一步提升数量推断的稳定性与准确性。

\par\noindent
（3）在模型压缩与部署验证方面，本文已从参数量、浮点运算量以及相关任务指标等角度对压缩方法进行了分析，但尚未在实际边缘硬件平台上完成系统化部署测试。后续可结合具体边缘设备，对压缩模型的推理时延、资源占用与运行稳定性进行进一步验证，从而更全面地评估模型在实际仓储场景中的应用效果。

\par\noindent
（4）在方法拓展方面，本文当前主要依赖单帧图像进行检测、分割与数量估计。未来可进一步结合视频时序信息、多模态感知信息或仓储管理系统中的先验业务信息，构建更加稳定和智能的库存盘点方法，以满足更复杂场景下的实际应用需求。
